\documentclass[11pt]{article}
\usepackage[margin=1in]{geometry}
\usepackage[T1]{fontenc}
\usepackage{lmodern}
\usepackage{booktabs}
\usepackage{tabularx}
\usepackage{enumitem}
\usepackage{xcolor}
\usepackage{titlesec}
\usepackage[hidelinks]{hyperref}

\definecolor{navy}{HTML}{1F3A5F}
\titleformat{\section}{\large\bfseries\color{navy}}{\thesection.}{0.5em}{}
\setlist{itemsep=2pt, topsep=4pt}
\newcolumntype{Y}{>{\raggedright\arraybackslash}X}

\title{\textbf{Hybrid Quantum-Classical Pneumonia Detection}\\[4pt]
\large COSC 461/561 Project 2: Project Specifications}
\author{Hasset Admasu \and Nathnael Abraham \and Andres Iniguez}
\date{November 2026}

\begin{document}
\maketitle

\section{Research Question}
Can a variational quantum circuit, simulated with the QuantumRings SDK and used as the
classification head on top of a finetuned CNN, match a classical head with the same number
of trainable parameters for pneumonia detection on chest X-rays, and does it hold up better
when training data is scarce?

\section{Dataset}
\begin{itemize}
  \item \textbf{Source:} PneumoniaMNIST (MedMNIST v2), installed via \texttt{pip install medmnist}.
  \item \textbf{Size:} 5,856 pediatric chest X-rays at $128\times128$ pixels, grayscale.
  \item \textbf{Task:} binary classification, normal vs.\ pneumonia.
  \item \textbf{Split:} official train / validation / test split; the test set is held out for final evaluation only.
  \item \textbf{Known issue:} class imbalance (about 74\% pneumonia), handled with a class-weighted loss or balanced sampling.
\end{itemize}

\section{Exploratory Data Analysis}
\begin{itemize}
  \item Class distribution per split.
  \item Sample image grid for each class.
  \item Pixel intensity histograms by class.
  \item Duplicate and corrupted image check.
  \item t-SNE or UMAP of CNN embeddings colored by label.
\end{itemize}

\section{Preprocessing}
\begin{itemize}
  \item Resize to $128\times128$ and normalize pixel values.
  \item Light augmentation on training data only: small rotations, translations, and brightness/contrast
        jitter (no horizontal flips, since chest anatomy is asymmetric).
\end{itemize}

\section{Models}
\begin{center}
\begin{tabularx}{\textwidth}{@{}l Y p{3.6cm}@{}}
\toprule
\textbf{Role} & \textbf{Architecture} & \textbf{Trainable components} \\
\midrule
Baseline 1 & Frozen ImageNet ResNet-18 features + logistic regression & Linear head \\
Baseline 2 & Finetuned ResNet-18 + small MLP head, parameter-matched to the quantum head & CNN + MLP head \\
Main approach & Finetuned ResNet-18 $\rightarrow$ linear projection to $n$ features $\rightarrow$ angle encoding
  into $n$ qubits $\rightarrow$ variational circuit (RY/RZ rotations + CNOT ring, $L$ layers)
  $\rightarrow$ Pauli-Z expectation $\rightarrow$ sigmoid & CNN, projection layer, circuit rotation angles \\
\bottomrule
\end{tabularx}
\end{center}

\noindent\textbf{Training procedure:} (1) finetune the CNN on the course GPU with a classical head;
(2) freeze the CNN and cache embeddings; (3) train the projection layer and quantum circuit on the
cached embeddings, with circuit gradients from the parameter-shift rule.

\medskip
\noindent\textbf{Quantum tooling:} QuantumRings SDK (\texttt{QuantumRingsLib} and/or
\texttt{quantumrings-toolkit-qiskit}), circuits of 4 to 8 qubits. The API token is stored in an
environment variable, never committed to the repository.

\section{Controlled Experiments}
\begin{center}
\begin{tabularx}{\textwidth}{@{}p{3.2cm} p{2.8cm} Y p{3.3cm}@{}}
\toprule
\textbf{Experiment} & \textbf{Variable} & \textbf{Values} & \textbf{Held constant} \\
\midrule
Main comparison & Head type & Logistic regression, MLP, quantum & Data, CNN backbone, splits \\
Change 1: data size & Fraction of training set & 1\%, 5\%, 25\%, 100\% & Test set, architecture \\
Change 2: circuit capacity & Qubits $\times$ layers & Qubits $\{4,6,8\}$ $\times$ layers $\{1,2,4\}$ & MLP resized to equal parameter count \\
Stretch: noise & Depolarizing noise rate & 0, 0.001, 0.01 & Best circuit configuration \\
\bottomrule
\end{tabularx}
\end{center}
Every setting is run with 3 random seeds; results report the mean and standard deviation.

\section{Evaluation}
\begin{itemize}
  \item \textbf{Primary metric:} AUROC (robust to class imbalance).
  \item \textbf{Secondary metrics:} F1 score, sensitivity (recall on pneumonia), specificity, accuracy, training time.
  \item \textbf{Figures:} learning curves vs.\ training data size, AUROC heatmap over qubits $\times$ layers,
        ROC curves, confusion matrices, and results tables per experiment.
\end{itemize}

\section{Engineering Requirements}
\begin{itemize}
  \item \textbf{Configuration:} Hydra or YAML configs; each run is launched by one command, e.g.\
        \texttt{python train.py head=quantum qubits=6 data\_frac=0.05 seed=1}.
  \item \textbf{Experiment tracking:} MLflow or Weights \& Biases logging parameters, metrics, and plots for every run.
  \item \textbf{Reproducibility:} fixed seeds, pinned dependencies, Makefile targets (\texttt{data}, \texttt{train}, \texttt{eval}).
  \item \textbf{Repository layout:} \texttt{src/data}, \texttt{src/models}, \texttt{src/train.py},
        \texttt{src/evaluate.py}, \texttt{configs/}, \texttt{notebooks/eda.ipynb}, \texttt{tests/}.
  \item \textbf{Testing and CI:} pytest checks (circuit output shape, equal parameter counts across heads,
        no train/test leakage) run by GitHub Actions on each push.
  \item \textbf{Demo:} FastAPI or Gradio inference endpoint packaged with a Dockerfile.
  \item \textbf{Documentation:} README with setup and run instructions, plus a model card stating intended
        use and that the model is not a clinical tool.
\end{itemize}

\section{Deliverables}
\begin{itemize}
  \item Progress presentation (10 minutes), due November 30, 5:30 PM.
  \item Final 5 to 7 page PDF report and GitHub repository, due December 11, 11:59 PM.
\end{itemize}

\end{document}
